This study compares interpretable and nonlinear classifiers for identifying existing gallstone status from demographic, laboratory, and bioimpedance measurements. A secondary analysis of 319 records with 38 candidate predictors evaluated logistic regression, random forest, XGBoost, LightGBM, and stacking. Training-fitted preprocessing and CART-based feature selection were assessed using five stratified training--validation--test splits and complementary leave-one-out cross-validation. Logistic regression achieved the highest mean validation ROC--AUC with direct CART selection (0.834) and recursive feature elimination (0.838), and was selected in all five fixed-model comparisons. The validation-selected procedure achieved mean internal test AUC 0.883, sensitivity 0.892, and specificity 0.708. Under leave-one-out evaluation, logistic regression had the highest AUC (0.874), whereas random forest had the highest sensitivity (0.839). Reducing the recursive-selection feature count from 20 to 10 lowered mean validation AUC from 0.838 to 0.816. These findings favour logistic regression as a parsimonious research baseline, while revealing sensitivity and specificity trade-offs. Overlapping and reused assessment partitions, a single clinical source, and the absence of calibration and external validation limit inference. The results concern existing disease status and do not establish prospective prediction or clinical utility.
